% !TeX program = tectonic
\documentclass[11pt,a4paper]{article}
\usepackage{fontspec}
\usepackage[russian]{babel}
\babelfont{rm}[Language=Default]{Liberation Serif}
\babelfont[Language=Default]{sf}{Carlito}
\babelfont[Language=Default]{tt}{DejaVu Sans Mono}
\usepackage{amsmath,amssymb,amsthm}
\usepackage{geometry}
\geometry{a4paper, top=2.4cm, bottom=2.4cm, left=2.4cm, right=2.4cm}
\usepackage{booktabs,tabularx,enumitem,xcolor,tcolorbox}
\tcbuselibrary{breakable,skins}
\definecolor{linkblue}{HTML}{1F4E79}
\definecolor{statgreen}{HTML}{2F6B3A}
\definecolor{goldacc}{HTML}{8B7E5A}
\usepackage[colorlinks=true,linkcolor=linkblue,urlcolor=linkblue,unicode=true]{hyperref}
\theoremstyle{plain}
\newtheorem{theorem}{Теорема}
\newtheorem{lemma}[theorem]{Лемма}
\newtheorem{proposition}[theorem]{Предложение}
\newtheorem{corollary}[theorem]{Следствие}
\theoremstyle{definition}
\newtheorem{definition}[theorem]{Определение}
\theoremstyle{remark}
\newtheorem{remark}[theorem]{Замечание}
\newtcolorbox{statusbox}[1][]{colback=statgreen!5,colframe=statgreen!75,
  title={\textbf{Сводка доказанного:} #1},fonttitle=\small,boxrule=0.7pt,arc=2pt,breakable}
\newtcolorbox{databox}[1][]{colback=goldacc!6,colframe=goldacc!80,
  title={\textbf{Данные протокола:} #1},fonttitle=\small,boxrule=0.7pt,arc=2pt,breakable}
\newtcolorbox{verifybox}[1][]{colback=linkblue!5,colframe=linkblue!80,
  title={\textbf{Верификация:} #1},fonttitle=\small,boxrule=0.7pt,arc=2pt,breakable}
\widowpenalty=10000
\clubpenalty=10000
\setlength{\parskip}{0.35em}
\newcommand{\CC}{\mathbb{C}}
\newcommand{\RR}{\mathbb{R}}
\newcommand{\ZZ}{\mathbb{Z}}
\newcommand{\QQ}{\mathbb{Q}}
\newcommand{\Ga}{\Gamma}
\newcommand{\im}{\operatorname{Im}}
\newcommand{\re}{\operatorname{Re}}
\newcommand{\disc}{\operatorname{disc}}
\begin{document}
\begin{center}
{\LARGE\bfseries Сертификат I: ступень Гурвица $N=7$}\\[0.4em]
{\large Циклическая кубика $2\cos(2\pi/7)$, спаривание Кардано и фаза $\pi/7$}\\[0.6em]
{\normalsize Исаев Исхак Хамзатович}\\[0.2em]
{\small Программа <<Динамический принцип>> --- лаборатория \texttt{hodge-laboratory}}\\[0.15em]
{\small Отдельная монография теоремы · Монография 17/20}
\end{center}
\vspace{0.4em}
{\color{linkblue}\hrule height 0.8pt}
\vspace{0.8em}

\section{Постановка}

С расширением программы до четырёх ступеней $N=7,9,15,30$ (дорожная
карта v1.3) лестница получила два кубических основания: ступень
Гурвица $N=7$ и ступень Макбета $N=9$. Сертификат I --- первая из
двух новых опор: он закрепляет на уровне $N=7$ весь конвейер
программы --- перепись характеров по проводникам, точный
кубический слой $b_{\mathrm{Ch}}(7)$ в $\ZZ[\zeta]/(\Phi_7)$ и
форму Кардано в $\mathrm{casus\ irreducibilis}$.

Имя ступени носит Адольф Гурвиц: многочлен $x^3+x^2-2x-1$ ---
минимальный многочлен $2\cos(2\pi/7)$, порождающего элемент
наибольшего вполне вещественного подполя кругового поля
$\QQ(\zeta_7)$ --- того самого поля, которое живёт в дискриминанте
$\Delta=-343=-7^3$ квартики Клейна (монография 16/20). Отметим,
однако, что стенд $N=7$ сертифицирует не саму квартику Клейна
(для неё есть отдельный стенд <<Клейн>> и монография 16/20), а
ферма-уровень конвейера при $N=7$: кривая Ферма $X_7$ имеет род
$15$, и все механизмы программы --- перепись, периоды,
$\Ga$-мономиалы --- работают на этом уровне без изменений.

\begin{theorem}[I1 --- перепись простого уровня]\label{th:I1}
Для кривой Ферма $X_7$ каждый собственный характер $(a,b)$,
$1\le a,b$, $a+b\le 6$, имеет проводник $d=7$; перепись
характеров по проводникам состоит из одного класса:
$\{7\colon 15\}$, и $\sum_d h_d = 15 = (7-1)(7-2)/2 = g$.
Обе схемы переписи (прямая и моёбиусова) дают один и тот же
гистограмм.
\end{theorem}

\begin{proof}
Уровень $N=7$ прост, поэтому для любого собственного характера
$(a,b)$ с $1\le a,b$, $a+b\le 6$ имеем $a,b\not\equiv 0 \pmod 7$,
откуда $\gcd(7,a,b)=1$ и проводник $d = 7/\gcd(7,a,b) = 7$ ---
единственный. Число таких пар $(a,b)$ --- точки целочисленного
треугольника $1\le a,b$, $a+b\le 6$: для каждого $a$ от $1$ до
$5$ ровно $6-a$ значений $b$, итого
$\sum_{a=1}^{5}(6-a)=5+4+3+2+1=15$. Моёбиусова схема:
$h_7 = c(7)-c(1) = 15-0 = 15$, где
$c(k)=(k-1)(k-2)/2$ --- число собственных характеров уровня $k$.
Обе схемы --- целочисленные тождества, их согласие
детерминировано.
\end{proof}

\begin{theorem}[I2 --- кубический слой $b_{\mathrm{Ch}}(7)$]\label{th:I2}
Числа $x_1=\zeta+\zeta^{-1}$, $x_2=\zeta^2+\zeta^{-2}$,
$x_3=\zeta^3+\zeta^{-3}$, $\zeta=e^{2\pi i/7}$, имеют точные
элементарные симметрические многочлены
\[
  s_1=x_1+x_2+x_3=-1,\qquad
  s_2=x_1x_2+x_1x_3+x_2x_3=-2,\qquad
  s_3=x_1x_2x_3=1,
\]
вычисленные чисто целочисленно в $\ZZ[\zeta]/(\Phi_7)$; поэтому
$x_1$ --- корень кубики $x^3+x^2-2x-1$, и это --- его точный
минимальный многочлен над $\QQ$.
\end{theorem}

\begin{proof}
\emph{Сумма.} $s_1=\zeta+\zeta^6+\zeta^2+\zeta^5+\zeta^3+\zeta^4$
--- сумма всех нетривиальных степеней $\zeta$, то есть
$s_1=\Phi_7$-сумма $= -1$ (так как $1+\zeta+\dots+\zeta^6=0$).

\emph{Попарные произведения.} Раскроем в $\ZZ[\zeta]/(\Phi_7)$,
пользуясь $\zeta^{-k}=\zeta^{7-k}$:
$x_1x_2=(\zeta+\zeta^6)(\zeta^2+\zeta^5)=\zeta^3+\zeta^6+\zeta+\zeta^4$;
$x_1x_3=(\zeta+\zeta^6)(\zeta^3+\zeta^4)=\zeta^4+\zeta^5+\zeta^2+\zeta^3$;
$x_2x_3=(\zeta^2+\zeta^5)(\zeta^3+\zeta^4)=\zeta^5+\zeta^6+\zeta+\zeta^2$.
Складывая, получаем каждую из шести степеней $\zeta^1,\dots,\zeta^6$
ровно дважды:
$s_2=2(\zeta+\zeta^2+\dots+\zeta^6)=2\cdot(-1)=-2$.

\emph{Произведение.}
$x_1x_2x_3=(\zeta^3+\zeta^6+\zeta+\zeta^4)(\zeta^3+\zeta^4)
=\zeta^6+\zeta^7+\zeta^9+\zeta^{10}+\zeta^4+\zeta^5+\zeta^7+\zeta^8
=\zeta^6+1+\zeta^2+\zeta^3+\zeta^4+\zeta^5+1+\zeta
=2+(\zeta+\dots+\zeta^6)=2-1=1$.

По формулам Виета кубика с корнями $x_1,x_2,x_3$ есть
$x^3-s_1x^2+s_2x-s_3=x^3+x^2-2x-1$. Минимальность: группа Галуа
$\mathrm{Gal}(\QQ(\zeta_7)^+/\QQ)\cong(\ZZ/7)^\times/\{\pm1\}\cong C_3$
действует умножением показателей на $2$:
$1\mapsto2\mapsto4\equiv-3\mapsto-1\equiv1$ --- цикл на трёх
классах $\{1,6\},\{2,5\},\{3,4\}$, то есть действие транзитивно и
$[\QQ(x_1):\QQ]=3$. Отпечаток: по модулю $2$ кубика
$x^3+x^2-2x-1\equiv x^3+x^2+1$ не имеет корней в $\mathbf{F}_2$
($0\mapsto1$, $1\mapsto1$) --- неприводим, независимо
подтверждая степень $3$.
\end{proof}

\begin{theorem}[I3 --- форма Кардано и спаривание]\label{th:I3}
В терминах $t=\frac{7}{54}+\frac{7\sqrt{-3}}{18}$ и
$\bar t=\frac{7}{54}-\frac{7\sqrt{-3}}{18}$ наибольший корень
кубики выражается как
\[
  2\cos\frac{2\pi}{7}=\sqrt[3]{t}+\sqrt[3]{\bar t}-\frac13,
  \qquad
  b_{\mathrm{Ch}}(7)=1-\cos\frac{2\pi}{7}
  =\frac76-\frac{\sqrt[3]{t}+\sqrt[3]{\bar t}}{2},
\]
причём модули аргументов Кардано спарены тождеством
$|t|^2=(7/9)^3$, так что произведение кубических корней равно
$\sqrt[3]{t}\cdot\sqrt[3]{\bar t}=7/9$ точно.
\end{theorem}

\begin{proof}
Подстановка $x=y-\frac13$ убирает квадратичный член:
$p=-\frac73$, $q=-\frac{7}{27}$, и кубика превращается в
$y^3-\frac73y-\frac7{27}=0$. Дискриминант Кардано
\[
  \Delta=\frac{q^2}{4}+\frac{p^3}{27}
  =\frac{49}{2916}-\frac{343}{729}
  =\frac{49-1372}{2916}=-\frac{1323}{2916}=-\frac{49}{108}<0
\]
--- классический \emph{casus irreducibilis}: все три корня
вещественны, но выражаются только через комплексные кубические
корни. Тогда $\sqrt{\Delta}=\frac{7i\sqrt3}{18}$ и
\[
  y=\sqrt[3]{-\frac q2+\sqrt\Delta}+\sqrt[3]{-\frac q2-\sqrt\Delta}
  =\sqrt[3]{\frac{7}{54}+\frac{7\sqrt{-3}}{18}}
  +\sqrt[3]{\frac{7}{54}-\frac{7\sqrt{-3}}{18}},
\]
откуда $x_1=y-\frac13$ и $b_{\mathrm{Ch}}(7)=1-\frac{x_1}{2}
=\frac76-\frac{\sqrt[3]{t}+\sqrt[3]{\bar t}}{2}$.

\emph{Спаривание.} $|t|^2=\frac{49}{2916}+\frac{147}{324}
=\frac{49+1323}{2916}=\frac{1372}{2916}=\frac{343}{729}=\left(\frac79\right)^3$,
поэтому $\sqrt[3]{t}\cdot\sqrt[3]{\bar t}=\sqrt[3]{|t|^2}
=\sqrt[3]{(7/9)^3}=7/9$ --- точное целочисленное тождество
(в коде лаборатории комментарий <<the Cardano pairing>>).

\emph{Виет-форма.} $|t|^{1/3}=\sqrt{7}/3$,
$\arg t=\arctan(3\sqrt3)$, поэтому главный кубический корень
$u=\sqrt[3]{t}=\frac{\sqrt7}{3}e^{i\theta}$,
$\theta=\frac13\arctan(3\sqrt3)$, и
$2\cos\frac{2\pi}{7}=2\re u-\frac13
=\frac{2\sqrt7}{3}\cos\!\Big(\frac13\arctan(3\sqrt3)\Big)-\frac13$
--- вещественная тригонометрическая форма (трисекция угла
$\arctan(3\sqrt3)\approx79.1^\circ$).

\emph{Изоляция.} Корни кубики $x_1\approx2.24698$,
$x_2\approx-0.55496$, $x_3\approx-1.69202$ строго упорядочены;
численная проверка при рабочей точности (dps 60) даёт
относительное расхождение формы Кардано с $1-\cos(2\pi/7)$
равным $2.1\cdot10^{-61}$ --- совпадение в пределах рабочей
точности, фиксирующее выбор старшего корня.
\end{proof}

\begin{remark}
Фаза $\pi/7$ --- якорь квантовой лестницы программы: пункт G6
сертификата G закрепляет фазы $\pi/7,\pi/15,\pi/30$ как уровни
одной конструкции $\delta=\pi/N$, а изотипия $(1,1,0,1)$ квартики
Клейна связывает $\zeta_7$ с её периодами. Ступень $N=7$ --- таким
образом, арифметическое основание всей гептадной ветви.
\end{remark}

\begin{databox}[Сертификат I --- числа стенда]
Род $g=15$; перепись $\{7\colon 15\}$ (обе схемы); минимальная
кубика $x^3+x^2-2x-1$; Виета $(s_1,s_2,s_3)=(-1,-2,1)$ точно в
$\ZZ[\zeta]/(\Phi_7)$; кубика неприводима над $\mathrm{GF}(2)$;
форма Кардано
$b_{\mathrm{Ch}}(7)=\frac76-\frac{\sqrt[3]{t}+\sqrt[3]{\bar t}}{2}$,
$t=\frac{7}{54}\pm\frac{7\sqrt{-3}}{18}$; спаривание
$\sqrt[3]{t}\sqrt[3]{\bar t}=7/9$;
$b_{\mathrm{Ch}}(7)=0.376510198141266469474995115995760\ldots$;
относительные ошибки протокола: Кардано $2.0\cdot10^{-36}$,
периоды (выборка) $1.0\cdot10^{-36}$ (dps 35); фаза $\pi/7$.
\end{databox}

\begin{statusbox}[итог]
Сертификат I принят ($5/5$): перепись простого уровня двумя
схемами, точные целые Виета в $\ZZ[\zeta]/(\Phi_7)$,
неприводимость над $\mathrm{GF}(2)$, форма Кардано в
$\mathrm{casus\ irreducibilis}$ с точным спариванием $7/9$,
фаза $\pi/7$. Кубическое основание лестницы установлено.
\end{statusbox}

\begin{verifybox}[как воспроизвести]
\texttt{python3 laboratory.py} $\to$ <<Сертификаты $\to$ I>>;
<<Стенды $\to$ N=7>>; batch-режим: сценарий с пунктом
\texttt{bch} ($N=7$); \texttt{--check-baseline} сверяет
$g=15$, симметрические $(-1,-2,1)$ и ошибки с эталоном.
\end{verifybox}

\vfill
{\color{linkblue}\hrule height 0.6pt}
\vspace{0.5em}
{\small\color{gray}
Лицензия: индивидуальная исключительная лицензия Исаева Исхака Хамзатовича.
Все права на вычисления, формулы и тексты принадлежат автору.
Верификация: \texttt{python3 laboratory.py} (пункты <<Стенды>>/<<Сертификаты>>).
}
\end{document}
